% Lösungen Einheit 4 – Prozentuale Veränderung
\einheitenkopf[][4 Veränderung]{Einheit 4 von 4 · Prozentuale Veränderung}

\erg{42}{a) auf $80\,\%$ gesenkt \quad b) auf $110\,\%$ erhöht \quad c) um $25\,\%$ herabgesetzt \quad d) um $25\,\%$ erhöht \quad e) um $40\,\%$ gesenkt}
\erg{43}{a) $40\,€$ \quad b) $10\,€$ \quad c) 60 Minuten \quad d) 1000 Besucher \quad e) $300\,€$}
\erg{44}{a) $6\,€$ dazu: $66\,€$ \quad b) $72\,€$ \quad c) $90\,€$ \quad d) $75\,€$ \quad e) Der alte Preis $60\,€$ bleibt; je größer der Prozentsatz, desto mehr kommt dazu.}
\erg{45}{a) $4\,€$ weg: $36\,€$ \quad b) $27\,€$ weg: $63\,€$}
\erg{46}{a) $40\,€ \cdot 1{,}2 = 48\,€$ \quad b) $70\,€ \cdot 0{,}8 = 56\,€$ \quad c) Faktor $1{,}05$; $200\,€ \cdot 1{,}05 = 210\,€$}
\erg{47}{a) $10 : 50 = 0{,}2 = 20\,\%$ mehr \quad b) $20 : 80 = 0{,}25 = 25\,\%$ weniger \quad c) Unterschied $0{,}07\,€$; $0{,}07 : 1{,}15 = 0{,}0608\ldots \approx 6{,}1\,\%$ (nicht $0{,}07 : 1{,}22 \approx 5{,}7\,\%$)}
\erg{48}{a) $80\,\% = 32\,€$, $1\,\% = 0{,}40\,€$, $100\,\% = 40\,€$ \quad b) $125\,\% = 10\,€$, $1\,\% = 0{,}08\,€$, $100\,\% = 8\,€$}
\erg{49}{a) $200\,€ + 38\,€ = 238\,€$ \quad b) $30\,€ + 2{,}10\,€ = 32{,}10\,€$ \quad c) $119\,\% = 59{,}50\,€$, $1\,\% = 0{,}50\,€$, $100\,\% = 50\,€$}
\erg{50}{a) 10 Prozentpunkte \quad b) $10 : 40 = 0{,}25 = 25\,\%$ \quad c) $5\,\%$ von 600 $= 30$ Kinder (nicht 5)}
\erg{51}{a) $12\,$m \quad b) $3 : 60 = 0{,}05 = 5\,\%$ \quad c) $13 : 280 = 0{,}0464\ldots \approx 4{,}6\,\%$; $4{,}6\,\% > 4\,\%$, der Weg ist nicht erlaubt.}
\erg{52}{a) Unterschied auf den neuen Preis bezogen; richtig $5 : 20 = 0{,}25 = 25\,\%$ \quad b) Unterschied auf den neuen Preis bezogen; richtig $10 : 50 = 0{,}2 = 20\,\%$}
\erg{53}{a) Der alte Preis ist $100\,\%$, weil die Erhöhung von ihm aus gerechnet wird. \quad b) Nein. Der Unterschied ist beide Male $10\,€$, aber einmal ist $40\,€$ das Ganze, einmal $50\,€$; von $40\,€$ ist $10\,€$ ein größerer Anteil ($25\,\%$ gegen $20\,\%$).}
\erg{54}{a) $100\,€ \cdot 0{,}8 = 80\,€$; $80\,€ \cdot 0{,}9 = 72\,€$ \quad b) $70\,€$ \quad c) Laden B; die $10\,\%$ bei A werden vom kleineren Preis $80\,€$ gerechnet und sind nur $8\,€$.}
